\documentclass{article}
\usepackage{tcee-repaso}
\tcetema[El euro, el SME y el BCE]{3.B.43}{El origen del euro: funcionamiento y evolución del Sistema Monetario Europeo. Los criterios de convergencia nominal. El Sistema Europeo de Bancos Centrales: objetivos e instrumentos. La política monetaria en la Eurozona desde 2009.}
\tcefecha{19/11/2026}
\begin{document}
\cabecerarepaso
\begin{columnas}

% ==========================================================================
\bloque[2,5]{Introducción}
\begin{itemize}
\item \textbf{Enganche}: «Cría fama y échate a dormir».
  \begin{itemize}
  \item \autor{Draghi} (26-7-2012): «lo que sea necesario… y créanme, será suficiente»; las primas bajaron sin comprar un bono.
  \item Hoy: euro en 21 Estados; depósito al 2,50\,\% desde el 16-9-2026 (energía).
  \end{itemize}
\item \textbf{Relevancia}: integración monetaria sin Tesoro común; 2.ª moneda ($\approx$20\,\% de las reservas).
  \begin{itemize}
  \item España: política monetaria desde Fráncfort desde 1999; auge y crisis \vertema{4.A.28}.
  \end{itemize}
\item \textbf{Contextualización}: \concepto{euro} desde 1-1-1999 (11 Estados); serpiente (1972) y SME (1979); Delors (1989), Maastricht (1992).
  \begin{itemize}
  \item Teoría: áreas monetarias óptimas (\autor{Mundell}) \vertema{3.B.16}; unión bancaria \vertema{3.B.44}; gobernanza \vertema{3.B.45}.
  \end{itemize}
\item \textbf{Problemática}: ¿por qué fracasó el SME? ¿Qué se exige para entrar? ¿Quién decide y con qué? ¿Sirve una política para 21?
\item \textbf{Estructura}: \hyperref[rep:I]{I} origen; \hyperref[rep:II]{II} SEBC; \hyperref[rep:III]{III} política desde 2009.
\end{itemize}
\esquemacante{%
\apartadocante{Introducción}{2,5}
\apartadocante{I. Del SME al euro}{8}
\subcante{1. De Bretton Woods a la serpiente}{1}
\subcante{2. El SME (1979-1998)}{4}
\subcante{3. Delors, Maastricht y las 3 fases}{1,5}
\subcante{4. Los criterios de convergencia}{1,5}
\apartadocante{II. El SEBC}{9,5}
\subcante{1. El marco institucional}{3}
\subcante{2. El objetivo y la estrategia}{2}
\subcante{3. Los instrumentos}{2,5}
\subcante{4. El papel internacional del euro}{2}
\apartadocante{III. La política monetaria desde 2009}{9}
\subcante{1. La crisis de la deuda}{4}
\subcante{2. Deflación y expansión cuantitativa}{1,5}
\subcante{3. Pandemia, inflación y normalización}{1,5}
\subcante{4. La talla única}{2}
\apartadocante{Conclusión}{1}}

% ==========================================================================
\bloque[8]{I. Del Sistema Monetario Europeo al euro\label{rep:I}}
\begin{ideaclave}
Con capitales libres, los tipos fijos pero ajustables eran asimétricos (los fijaba el Bundesbank) y frágiles (1992-1993): la salida fue una moneda y un banco central comunes, con una entrada condicionada a la convergencia nominal.
\end{ideaclave}

\sub{I.1. De Bretton Woods a la serpiente\label{rep:I.1}}
\begin{itemize}
\item Plan Barre (1969), La Haya (1969), \emph{Informe Werner} (1970): UEM en 3 etapas antes de 1980.
  \begin{itemize}
  \item Lo frenan la crisis del dólar y la de 1973.
  \end{itemize}
\item Smithsonian (1971): $\pm2{,}25$\,\% frente al dólar (túnel de 4,5\,\%).
\item \concepto{Serpiente} (Basilea, 1972): máximo 2,25\,\% entre monedas europeas; desde 1973, «en el lago»; FECOM (1973).
  \begin{itemize}
  \item Salen la libra (1972), la lira (1973) y el franco (1974 y 1976): faltaban coordinación y apoyo.
  \end{itemize}
\end{itemize}
\graficorepaso[0.85\linewidth]{serpiente}{La serpiente en el túnel (1972-1973) y en el lago}

\sub{I.2. El Sistema Monetario Europeo (1979-1998)\label{rep:I.2}}
\subsub{a) Origen, objetivos y funcionamiento}
\begin{itemize}
\item \autor{Giscard} y \autor{Schmidt}; 13-3-1979; tipos \emph{fijos pero ajustables} con realineamientos.
\item \concepto{ECU}: cesta; unidad de cuenta.
\item \concepto{MTC}: $\pm2{,}25$\,\% ($\pm6$\,\%: lira, peseta desde 1989, libra 1990-1992).
  \begin{itemize}
  \item Intervención obligatoria e ilimitada en el límite; intramarginal con el indicador de divergencia (75\,\%).
  \item FECOM: créditos.
  \end{itemize}
\item \emph{Etapas}: inestable (1979-1983); convergencia (giro francés de 1983); SME «duro» (Basilea-Nyborg, 1987); crisis (1992-1993); bandas del $\pm15$\,\% (2-8-1993).
\end{itemize}
\subsub{b) La asimetría del SME: el problema $n-1$}
\begin{itemize}
\item \emph{Idea}: $n$ monedas, $n-1$ tipos de cambio: alguien fija la política; en el SME, el Bundesbank.
\item \emph{Supuestos}: ancla A y país B; tipo fijo, capitales libres; precios rígidos; el débil pierde reservas, el fuerte esteriliza.
\item \emph{Desarrollo}: A fija $M^S_A$ y con ello $i_A$; la oferta de B se vuelve endógena:
\funcionclave{i_B=i_A+\delta^e}
  \begin{itemize}
  \item Si A contrae (reunificación), B sube tipos aunque esté en recesión; si dudan de la paridad, $\delta^e\uparrow$.
  \end{itemize}
\end{itemize}
\graficorepaso[0.92\linewidth]{sme-asimetria}{El ancla fija $i_A$; B la sigue y pierde reservas}
\begin{itemize}
\item \emph{Implicaciones}: credibilidad importada («atarse las manos», \autor{Giavazzi}-\autor{Pagano}, 1988) \vertema{3.B.16}; coste con perturbaciones asimétricas.
\item \emph{Extensiones}: ataque autocumplido y cláusula de escape (\autor{Obstfeld}, 1996) \vertema{3.B.17}.
\item \emph{Evidencia}: los tipos alemanes guían (\autor{Giavazzi}-\autor{Giovannini}, 1989); hegemonía incompleta (\autor{Fratianni}-\autor{Von Hagen}, 1990).
\item \emph{Valoración}: Francia e Italia querían voz: banco central común.
\end{itemize}
\subsub{c) La crisis de 1992-1993}
\begin{itemize}
\item Causas: 5 años sin realinear; reunificación y Bundesbank; capitales libres; «no» danés y «sí» francés justo.
\item 16-9-1992: sale la libra; el 17, la lira.
\item Peseta: $-5$\,\% (sept. 1992), $-6$\,\% (nov.), $-8$\,\% (mayo de 1993), $-7$\,\% (1995).
\item Lección: flotar o compartir moneda.
\end{itemize}

\sub{I.3. El camino a la UEM: Delors, Maastricht y las 3 fases\label{rep:I.3}}
\begin{itemize}
\item \emph{Por qué}: mercado único, salir de la asimetría, trinidad imposible; Alemania reunificada; voz para Francia.
\item \emph{Informe Delors} (1989); \emph{Maastricht} (7-2-1992; vigor 1-11-1993): BCE independiente, criterios, calendario; \emph{opt-out} de RU y Dinamarca.
\item \emph{Fases}:
  \begin{enumerate}
  \item 1-7-1990: capitales libres (Ley 13/1994 del BdE);
  \item 1-1-1994: IME; «euro» (Madrid, 1995); PEC (1997); mayo de 1998, 11 Estados; BCE el 1-6-1998;
  \item 1-1-1999: 1 € = 166,386 pesetas; billetes, 1-1-2002.
  \end{enumerate}
\item Grecia (2001) $\to$ Bulgaria (2026): 21.
\end{itemize}
\graficorepaso[0.95\linewidth]{cronologia-euro}{Del Informe Werner al euro de 21}

\sub{I.4. Los criterios de convergencia nominal\label{rep:I.4}}
\begin{itemize}
\item Art.~140 TFUE y Protocolo 13:
  \begin{itemize}
  \item \emph{inflación} $\leq$ 3 mejores $+1{,}5$ pp;
  \item \emph{finanzas}: sin procedimiento de déficit excesivo (3\,\% y 60\,\%);
  \item \emph{tipo de cambio}: 2 años en el MTC II sin tensiones ($\pm15$\,\%; Dinamarca $\pm2{,}25$);
  \item \emph{tipo largo} $\leq$ 3 mejores $+2$ pp;
  \item convergencia legal (art.~131).
  \end{itemize}
\item Informes del BCE y de la Comisión cada 2 años.
  \begin{itemize}
  \item Bulgaria (2025) cumple: euro el 1-1-2026 (1,95583 leva).
  \item 2026: Chequia y Suecia cumplen precios, tipos y finanzas; nadie en el MTC II ni con convergencia legal.
  \end{itemize}
\item \emph{Valoración}: credibilidad (España: inflación del 6 a $<$2\,\%); pero nominales: nada sobre crédito ni cuenta corriente; endogeneidad (\autor{Frankel}-\autor{Rose}).
\end{itemize}
\graficorepaso[0.92\linewidth]{convergencia-2026}{Informe de Convergencia de 2026: inflación y tipos frente a la referencia}

% ==========================================================================
\bloque[9,5]{II. El Sistema Europeo de Bancos Centrales\label{rep:II}}
\begin{ideaclave}
Un banco central muy independiente, con mandato jerárquico (2\,\% simétrico a medio plazo) y un marco operativo que ha pasado del corredor al suelo con demanda.
\end{ideaclave}

\sub{II.1. El marco institucional\label{rep:II.1}}
\subsub{a) SEBC, Eurosistema y BCE}
\begin{itemize}
\item Arts.~127-133 y 282-284 TFUE; Protocolo 4.
\item \concepto{SEBC} (BCE + 27); \concepto{Eurosistema} (BCE + 21); BCE en Fráncfort, institución desde Lisboa.
\item Funciones: política monetaria, divisas, reservas, pagos (TARGET), billetes; supervisión desde 2014 \vertema{3.B.44}.
\item Exclusiva (art.~3.1.c); \emph{tipo de cambio} (art.~219): acuerdos u orientaciones, nunca usados; el euro flota.
\end{itemize}
\subsub{b) La independencia (art.~130)}
\begin{itemize}
\item Institucional; personal (8 años no renovables; gobernadores $\geq$5); funcional (art.~123: sin financiación monetaria; 124 y 125); financiera.
\item Rendición de cuentas ante el PE; cuentas de las reuniones (2015).
\end{itemize}
\subsub{c) Los órganos rectores}
\begin{itemize}
\item \concepto{Comité Ejecutivo} (6): \autor{Lagarde} hasta el 7-2-2027; \autor{Vujčić} VP desde 1-6-2026; \autor{Lane}, \autor{Schnabel}, \autor{Cipollone}, \autor{Elderson}.
\item \concepto{Consejo de Gobierno} (6 + 21): decide cada 6 semanas por mayoría simple.
  \begin{itemize}
  \item \emph{Rotación} (2015): grupo 1, los 5 mayores, 4 votos (80\,\% de los meses); grupo 2, los otros 16, 11 votos (69\,\%).
  \item 6 + 4 + 11 = 21 votos; con 22, 3 grupos.
  \end{itemize}
\item \concepto{Consejo General} (2 + 27): enlace con los de fuera; transitorio.
\item Clave de capital: 50\,\% población, 50\,\% PIB; cada 5 años.
\end{itemize}
\graficorepaso[0.92\linewidth]{organos-sebc}{SEBC, Eurosistema, BCE y órganos rectores}

\sub{II.2. El objetivo y la estrategia\label{rep:II.2}}
\begin{itemize}
\item Art.~127.1: estabilidad de precios y, sin perjuicio, apoyo a las políticas generales: \concepto{mandato jerárquico}.
\item \emph{Estrategias}:
  \begin{itemize}
  \item 1998: $<$2\,\% y 2 pilares (M3: 4,5\,\%); 2003: «inferior pero próxima al 2\,\%»;
  \item 2021: \concepto{2\,\% simétrico} a medio plazo; respuesta enérgica en el límite inferior; vivienda en propiedad; clima;
  \item 2025: respuesta a desviaciones grandes y duraderas en los 2 sentidos; escenarios; próxima, 2030.
  \end{itemize}
\item Por qué el 2\,\%: sesgo del IAPC, margen frente al límite inferior, ajustes relativos.
\end{itemize}

\sub{II.3. Los instrumentos\label{rep:II.3}}
\subsub{a) El marco operativo: del corredor al suelo}
\begin{itemize}
\item Variable operativa: \concepto{€STR} (2019; eonia hasta 3-1-2022).
\item \emph{Convencionales}: OPF (tipo fijo y sin límite desde 2008), LTRO a 3 meses; facilidades marginal (techo) y de depósito (suelo).
  \begin{itemize}
  \item Reservas mínimas 1\,\% (2012), sin remunerar desde 20-9-2023.
  \end{itemize}
\item \emph{Marco de 2024}: orientación con la facilidad de depósito; \concepto{suelo con demanda}; OPF a 15 pb.
\end{itemize}
\pareja{tipos-bce}{Tipos del BCE: 4\,\% (2023-2024), $-0{,}50$\,\% (2019-2022), 2,50\,\% hoy}{marco-operativo}{Corredor con reservas escasas; suelo con abundantes}
\subsub{b) Los instrumentos no convencionales}
\begin{itemize}
\item LTRO a 3 años, TLTRO (2014-2024); tipos negativos (2014-2022, mínimo $-0{,}50$\,\%).
\item APP (2015) y PEPP (2020; 1,85 billones); orientación sobre los tipos (2013).
\item SMP (2010); \concepto{OMT} (2012: con programa del MEDE; nunca usadas); \concepto{TPI} (2022: sin usar).
\item Balance: 8,8 billones (2022) $\to$ 5,9 (nov. 2026).
\end{itemize}
\subsub{c) El euro digital}
\begin{itemize}
\item Propuesta de 2023; Consejo (2025) y PE (2026); piloto 2027, emisión 2029; límite de tenencia, sin intereses.
\end{itemize}

\sub{II.4. El papel internacional del euro\label{rep:II.4}}
\begin{itemize}
\item Funciones del dinero $\times$ uso público/privado (\autor{Kenen}, 1983; \autor{Krugman}, 1984).
\item \concepto{Segunda moneda} (BCE, 2026): índice $\approx$20\,\%.
  \begin{itemize}
  \item Reservas 20,2\,\% (dólar 56,7); deuda 24,6; SWIFT $\approx$23; facturación de exportaciones 60\,\%.
  \end{itemize}
\item \emph{Depende de}: tamaño, estabilidad, \emph{activo seguro} y mercados profundos (punto débil), geopolítica.
  \begin{itemize}
  \item Gana señoreaje y financiación barata; arriesga apreciación en crisis.
  \end{itemize}
\item \emph{Apoyo}: \emph{swaps} y \concepto{EUREP} permanente y global (2026).
\item «Momento del euro global» (\autor{Lagarde}, 2025): unión de ahorros e inversiones, deuda segura, euro digital \vertema{3.B.42}.
\end{itemize}

% ==========================================================================
\bloque[9]{III. La política monetaria desde 2009\label{rep:III}}
\begin{ideaclave}
El BCE pasó a ser el respaldo de hecho de los Estados del euro: la credibilidad de que el euro es irreversible es su primer instrumento; pero solo puede hacer la política de la media.
\end{ideaclave}

\sub{III.1. La crisis de la deuda soberana (2009-2014)\label{rep:III.1}}
\subsub{a) Los hechos}
\begin{itemize}
\item Desequilibrios: tipos reales bajos en la periferia, crédito y vivienda; cuenta corriente (España $-9{,}6$\,\% en 2007).
\item Grecia (2009); SMP (2010); Deauville (2010); Irlanda y Portugal.
\item Subidas de \autor{Trichet} (2011), deshechas por \autor{Draghi}; LTRO a 3 años ($>$1 billón).
\item Diferenciales (julio de 2012): Grecia 2.739 pb, Portugal 1.203, España 555, Italia 519. 26-7-2012 y OMT.
\end{itemize}
\subsub{b) Las crisis de deuda autocumplidas}
\begin{itemize}
\item \emph{Idea} (\autor{Calvo}, 1988; \autor{De Grauwe}, 2011): deuda en una moneda que no se emite: el miedo al impago sube el tipo y hace el impago más probable.
\item \emph{Supuestos}: deuda $B$; superávit $S\sim U[0,\bar S]$; impago si $S<(1+i)B$; inversores neutrales (tipo $r$); sin banco central propio.
\item \emph{Desarrollo}: $p=(1+i)B/\bar S$; $1+i=(1+r)/(1-p^e)$; $\kappa=(1+r)B/\bar S$:
\funcionclave{p\,(1-p)=\kappa}
  \begin{itemize}
  \item $\kappa<\tfrac14$: 2 equilibrios, $p_L$ (estable) y $p_H$ (pérdida del acceso al mercado);
  \item $\kappa>\tfrac14$: insolvencia.
  \end{itemize}
\item \emph{Implicaciones}: el respaldo creíble elimina el equilibrio malo sin comprar; solo para solventes: condicionalidad.
\item \emph{Extensiones}: círculo banca-Estado \vertema{3.B.44}; redenominación.
\item \emph{Evidencia}: diferenciales sobre los fundamentos y caída tras las OMT (\autor{De Grauwe}-\autor{Ji}, 2013); $-2$ pp a 2 años en Italia y España (\autor{Altavilla} et al., 2016).
\item \emph{Valoración}: separar iliquidez e insolvencia es difícil: los tribunales.
\end{itemize}
\pareja{crisis-autocumplida}{$p=\kappa/(1-p^e)$ y la bisectriz: $p_L$, $p_H$ y el respaldo}{diferenciales}{Diferenciales a 10 años frente a Alemania}
\begin{itemize}
\item \emph{Tribunales}: \emph{Gauweiler} (2015): OMT legales; \emph{Weiss} (2018): PSPP legal; BVerfG (2020): \emph{ultra vires}.
\end{itemize}

\sub{III.2. Deflación y expansión cuantitativa (2014-2019)\label{rep:III.2}}
\begin{itemize}
\item Inflación negativa (enero de 2015: $-0{,}6$\,\%).
\item Tipos negativos (2014), TLTRO, APP (2015: 60.000 M€/mes; 80.000 en 2016-2017), orientación (\autor{Woodford}).
\end{itemize}

\sub{III.3. Pandemia, inflación y normalización (2020-2026)\label{rep:III.3}}
\begin{itemize}
\item \emph{Pandemia}: PEPP (18-3-2020; 1,85 billones); TLTRO III al $-1$\,\%.
\item \emph{Inflación}: 10,6\,\% (oct. 2022); +450 pb en 14 meses (depósito al 4\,\%); TPI (2022).
\item \emph{Normalización}: fin de reinversiones (APP 2023, PEPP 2024); 8 recortes (2024-2025) hasta el 2\,\%.
\item \emph{2026}: Oriente Próximo y energía: inflación 4,0\,\% (oct.; subyacente 2,5\,\%); subidas al 2,25\,\% (17-6) y al 2,50\,\% (16-9).
\end{itemize}
\pareja{balance-eurosistema}{Activo del Eurosistema: máximo de 8,8 billones}{inflacion-zona-euro}{IAPC general y subyacente}

\sub{III.4. Una política para 21 economías: la talla única\label{rep:III.4}}
\begin{itemize}
\item \emph{Idea}: el BCE mira la media; a cada país le conviene otro tipo.
\item \emph{Supuestos}: $i_j^*=r^*+\pi^*+\phi_\pi(\pi_j-\pi^*)+\phi_x x_j$ (\autor{Taylor}, 1993); el BCE, la misma regla con la media:
\funcionclave{i-i_j^*=\phi_\pi(\bar\pi-\pi_j)+\phi_x(\bar x-x_j)}
  \begin{itemize}
  \item Tipo real más bajo donde la inflación es más alta (\autor{Walters}); el canal de la competitividad, lento.
  \end{itemize}
\item \emph{Implicaciones}: 2002-2006, regla del 7-8\,\% para España y 2-4\,\% para Alemania, con el depósito al 1-2\,\%.
  \begin{itemize}
  \item Instrumentos nacionales: fiscal y macroprudencial.
  \end{itemize}
\item \emph{Extensiones}: fragmentación (TPI, OMT); sincronización endógena de los ciclos.
\item \emph{Evidencia}: inversión en 2012-2015; 2022, bálticos $>$15\,\%; 2026: España 7,5\,\%, Alemania 4,3\,\%.
\item \emph{Valoración}: la regla es solo una referencia; el BCE solo puede hacer la política de la media.
\end{itemize}
\graficorepaso[0.92\linewidth]{taylor-paises}{Regla de Taylor por país frente a la facilidad de depósito}

% ==========================================================================
\bloque[1]{Conclusión}
\begin{itemize}
\item \textbf{Recapitulación}: el euro nace de los fracasos cambiarios; Maastricht: convergencia nominal y BCE independiente; desde 2009, respaldo, compras y lucha contra la inflación.
\item \textbf{Relevancia}: para España, estabilidad y tipos bajos, pero también el auge del crédito y la lección de 2010-2012 \vertema{4.A.28}.
\item \textbf{Extensiones}: \vertema{3.B.16} \vertema{3.B.17} \vertema{3.B.44} \vertema{3.B.45} \vertema{4.A.19}.
\item \textbf{Opinión}: euro más sólido pero incompleto.
  \begin{itemize}
  \item Faltan el activo seguro, el seguro de depósitos, la unión de capitales y el euro digital.
  \end{itemize}
\item \textbf{Idea final}: la fama de 2012 bastó contra la ruptura; 2021 enseñó que no se puede dormir: la credibilidad se cría cumpliendo el 2\,\%.
\end{itemize}

% ==========================================================================
\anexorepaso
\bloque{Datos clave}
\begin{itemize}
\item Tipos: \dato{depósito 2,50\,\%}{16-9-2026}; máximo 4,00\,\% (2023-2024); mínimo $-0{,}50$\,\%.
\item Inflación: \dato{4,0\,\%}{oct. 2026}; máximo 10,6\,\% (oct. 2022).
\item Balance: \dato{5,9 billones}{nov. 2026}; máximo 8,8.
\item Banco de España: \dato{9,67\,\% del capital}{BCE}.
\end{itemize}
\bloque{Autores y fechas}
\hito{1961}{\autor{Mundell}: áreas monetarias}
\hito{1970}{Informe Werner}
\hito{1979}{SME}
\hito{1988}{\autor{Giavazzi}-\autor{Pagano}; \autor{Calvo}}
\hito{1992}{Maastricht; crisis del SME}
\hito{1993}{\autor{Taylor}: regla}
\hito{2011}{\autor{De Grauwe}: fragilidad}
\hito{2012}{\autor{Draghi}; OMT}
\bloque{Preguntas del tribunal}
\begin{enumerate}
\item \textbf{¿SEBC o Eurosistema?} BCE + 27 frente a BCE + 21.
\item \textbf{¿Por qué fracasó el SME?} Asimetría, reunificación, capitales libres, dudas sobre Maastricht.
\item \textbf{¿Compra el BCE deuda pública?} En el secundario, sí (\emph{Gauweiler}, \emph{Weiss}); en el primario, no (art.~123).
\item \textbf{¿OMT o TPI?} Con programa del MEDE frente a criterios de elegibilidad sin programa.
\item \textbf{¿Mandato dual?} No: jerárquico.
\item \textbf{¿Por qué criterios nominales y no reales?} Credibilidad; se esperaba la convergencia real (endogeneidad).
\end{enumerate}
\begin{trampa}
\begin{itemize}
\item Bandas del $\pm15$\,\% desde el \emph{2-8-1993}; la libra salió el \emph{16-9-1992}.
\item Inflación: frente a los 3 con \emph{mejor comportamiento}.
\item BCE: \emph{1-6-1998}; euro: 1-1-1999; billetes: 1-1-2002.
\item Serpiente: 2,25\,\% \emph{entre} monedas europeas; túnel: $\pm2{,}25$\,\% frente al dólar.
\item Reservas mínimas \emph{sin remunerar} desde el 20-9-2023.
\end{itemize}
\end{trampa}

\end{columnas}
\end{document}
